\documentclass[11pt]{article}
\usepackage[margin=1in]{geometry}
\usepackage{booktabs}
\usepackage{longtable}
\usepackage{amsmath}
\usepackage{graphicx}
\usepackage{float}
\usepackage{enumitem}
\usepackage[colorlinks=true,linkcolor=blue,urlcolor=blue]{hyperref}
\usepackage{xcolor}

\newcommand{\accepted}{\textcolor{teal}{\textbf{Accepted}}}
\newcommand{\preliminary}{\textcolor{orange}{\textbf{Preliminary}}}
\newcommand{\unresolved}{\textcolor{red}{\textbf{Unresolved}}}
\newcommand{\notrun}{\textcolor{gray}{\textbf{Not yet run}}}

\title{Repeating-Earthquake Coda-Wave Interferometry on SAFOD Borehole DAS\\
\large Progress log: what has been tried, what it returned}
\author{}
\date{\today}

\begin{document}
\maketitle
\tableofcontents
\newpage

\section{Why this project exists}

The active-source route was closed by measurement, not by opinion. The
accelerated weight drop (AWD) has an injection--recovery detection floor of about
$1\%$ apparent-velocity change. The largest published fault-zone $dv/v$ at
Parkfield outside an M6 mainshock is $0.02$--$0.2\%$: $0.08\%$ coseismic for the
2004 M6.0 (Brenguier et al.\ 2008, \texttt{10.1126/science.1160943}), at most
$0.2\%$ in the depth-resolved reanalysis (Wu et al.\ 2016,
\texttt{10.1002/2016GL069145}), and seasonal variation at Parkfield reported as
\emph{not evident} (Okubo et al.\ 2024, \texttt{10.1029/2023JB028084}).

A 5--50$\times$ gap set by source physics before acquisition cannot be closed by
processing. Prior art in the same borehole makes the point sharply: Niu et al.\
(2008, \texttt{10.1038/nature07111}) needed a 2\,kHz crosswell geometry over a
10\,m baseline to reach $0.0011\%$, and explicitly predicted the solid-earth tide
at SAFOD to be undetectable.

\subsection*{What replaces it}

Repeating earthquakes are a naturally repeatable source at seismogenic depth.
Coda-wave interferometry (CWI) on doublets routinely resolves
$\mathbf{0.02}$--$\mathbf{0.2\%}$ (Poupinet, Ellsworth \& Fr\'echet 1984,
\texttt{10.1029/JB089iB07p05719}) --- one to two orders of magnitude better than
first-arrival timing.

The reason is geometric. A coda wave has scattered many times and travelled a far
longer path than the direct arrival, so a small \emph{fractional} velocity change
accumulates into a large \emph{absolute} delay. Timing a runner over a marathon
rather than over ten metres.

\subsection*{Novelty}

No prior work anywhere combines repeating-earthquake CWI with DAS --- in a
borehole or otherwise. The two existing DAS+CWI studies use short cable segments
and ambient noise. Moreover, Lellouch et al.\ (2019,
\texttt{10.1029/2019JB017533}), working on \emph{this exact fiber}, proposed the
step and never took it: ``coda waves and converted modes may be examined more
systematically.'' Six years on, nobody has.

\subsection*{The framing, which matters more than the sensitivity argument}

Do \emph{not} pitch this as ``repeaters beat the weight drop.'' Rubinstein \&
Beroza (2005, \texttt{10.1029/2005GL023189}) used the \emph{SAFOD target
repeaters specifically} and concluded the coseismic velocity change is confined
to the top $<100$\,m, with direct S at depth much less affected. That is an
adverse, on-target precedent.

The framing that survives it is the one the geometry uniquely enables:
\textbf{depth-resolved $dv/v$ from a continuous vertical array}, turning their
surface-versus-100\,m binary into a continuous 0--919\,m profile and settling
where the change lives. This framing does not require the signal to be large.

\section{Data}

\begin{table}[H]
\centering
\begin{tabular}{p{0.35\linewidth}p{0.55\linewidth}}
\toprule
Quantity & Value \\
\midrule
Continuous DAS span & 2024-05-07 to 2025-09-24 \\
Covered days & \textbf{286} \\
Manifest rows (usable) & 392{,}038 \\
Fiber & 900 channels $\times$ 1.0210\,m = \textbf{919\,m} \\
Gauge length & 16.34\,m \\
Cable identity & Lellouch et al.\ (2019) cemented fiber \\
Catalog events in window & 540 \\
Events on covered days & \textbf{305 (56\%)} \\
\bottomrule
\end{tabular}
\caption{Continuous cemented-fiber archive and the Parkfield catalog over the
same window. Catalog from USGS FDSN, box $35.70$--$36.10^\circ$N,
$-120.70$--$-120.20^\circ$E, all magnitudes.}
\end{table}

The fiber spans 0--919\,m; Parkfield repeaters sit at 2--12\,km. The fiber does
not need to reach the source --- CWI needs a repeatable source--receiver pair, and
a 919\,m vertical array recording near-vertically-incident energy from below
gives depth-continuous sampling of the receiver-side path, which a surface
network cannot.

\section{A short primer on the methods used}

\paragraph{Coda-wave interferometry.} The coda is the scattered tail of the
seismogram. Because coda energy has travelled far longer paths than the direct
arrival, a uniform fractional velocity change $dv/v$ delays coda arriving at
lapse time $t$ by $\Delta t \approx -(dv/v)\,t$. Measuring that proportionality
is the measurement.

\paragraph{The stretching estimator.} Rather than picking delays, one resamples
(``stretches'') the later trace by a trial factor $\varepsilon$ and finds the
$\varepsilon$ maximising correlation with the earlier trace over a lapse window.
That $\varepsilon$ \emph{is} $dv/v$.

\paragraph{Why repeaters are required.} Coda looks like noise but is a
deterministic fingerprint of one source in one medium. Move the source tens of
metres and the fingerprint changes entirely. Repeating earthquakes --- the same
patch rupturing the same way --- are the only natural source that satisfies this.

\paragraph{Gates.} Each stage has a numeric pass/fail condition written before it
runs, so ``it produced a number'' cannot be confused with ``it worked.''

\section{What has been tried, and what it returned}

\subsection{Catalog intersection --- \accepted}

\textbf{Script:} \texttt{count\_doublets.py}. \textbf{Question:} are there enough
co-located, similar-magnitude events on days the fiber was recording to be worth
running waveform correlation on?

Result: 305 of 540 catalog events fall on covered days; 111 candidate co-located
pairs; \textbf{32 candidate sequences} with $\ge 2$ covered events; 13 within
5\,km of SAFOD; 30 at seismogenic depth.

\begin{table}[H]
\centering
\begin{tabular}{rrrrl}
\toprule
seq & $n$ & depth (km) & mag & span \\
\midrule
0 & 13 & 1.2--2.3 & 0.42--1.15 & 2024-05 $\to$ 2025-04 \\
2 & 12 & 3.5--4.3 & 0.04--1.14 & 2024-06 $\to$ 2025-04 \\
4 &  5 & 2.2--3.1 & \textbf{1.58--1.86} & 2024-06 $\to$ 2025-04 \\
5 &  4 & \textbf{9.87--10.05} & 0.95--1.45 & 2024-07 $\to$ 2025-03 \\
\bottomrule
\end{tabular}
\caption{Leading candidate sequences, all within $\sim$3--4\,km of SAFOD.}
\end{table}

\paragraph{Two defects in this list, found by inspecting it.} Single-link
clustering \emph{chains}: sequence 0 spans 1\,km of depth despite a 0.5\,km
threshold, so these are over-merged and the true sequence count is lower. And
several ``repeats'' are seconds apart (2025-03-29 06:43:47 / 06:43:58), which
provide no temporal baseline for CWI. Sequence 4 has the most repeater-like
magnitude tightness (0.28 spread); sequence 5 the tightest depths (180\,m).
Neither ranking survives without waveform confirmation.

\paragraph{Standing limitation.} Catalog absolute locations at Parkfield carry
$\sim$0.2--0.5\,km error while true repeaters are $<$20\,m across. Every number
here is an \textbf{upper bound} on candidates, not a repeater count.

\subsection{Gate G1 --- coda SNR --- \accepted (PASS)}

\textbf{Script:} \texttt{g1\_coda\_snr.py}. \textbf{Question:} does a 16.34\,m
gauge length give usable \emph{coda} SNR for M$\sim$1 events 3--5\,km away?
Nobody had demonstrated this. Lellouch showed \emph{arrivals}; coda is tens of dB
weaker and a much harder bar.

\textbf{Pass condition, written before running:} coda SNR $>3$ over $\ge 200$
channels for a 2--10\,s lapse window.

\begin{table}[H]
\centering
\begin{tabular}{lrrr}
\toprule
Event & Mag & Median coda SNR & Channels $>3$ \\
\midrule
2024-06-09 16:35:44 & 1.85 & 34.62 & 900/900 \\
2024-06-22 19:26:50 & 1.58 &  5.84 & 739/900 \\
2025-04-06 08:23:21 & 1.86 &  6.82 & 750/900 \\
2025-04-06 13:51:22 & 1.62 &  5.87 & 705/900 \\
2025-04-08 13:11:32 & 1.62 &  6.43 & 699/900 \\
\midrule
\multicolumn{3}{l}{Across events (median-SNR channels $>3$)} & \textbf{785/900} \\
\bottomrule
\end{tabular}
\caption{G1 result. \textbf{PASS} by roughly $4\times$ the threshold.}
\end{table}

Visual confirmation: the arrival is flat across all 900 channels (near-vertical
incidence), with clear coda in the lapse window.

\paragraph{An unplanned corroboration.} Median coda SNR drops below 3 in the
deepest $\sim$70\,m of fiber --- exactly where Lellouch et al.\ report the end-loop
failure at 800\,m. Two independent measurements, same boundary.

\subsection{Gate G2 --- waveform confirmation --- \unresolved{} (twice failed, rebuilt)}

\textbf{Script:} \texttt{g2\_confirm\_repeaters.py}. \textbf{Question:} are the
catalog clusters actually repeating earthquakes? \textbf{Pass:} $\ge 3$ sequences
with median cross-channel CC $>0.9$ for a pair separated by $>30$ days.

\paragraph{Attempt 1 --- FAIL, and the failure was mine.} Max off-diagonal CC of
$0.42$, $0.18$, $0.26$ across five sequences. The estimator computed
\emph{zero-lag} correlation. Catalog origin times for M$\sim$1 events carry
$0.1$--$0.5$\,s uncertainty, which at 100\,Hz is 10--50 samples; two
\emph{identical} waveforms offset by that much give near-zero zero-lag
correlation. Repeater identification is always done by maximising over lag.

\paragraph{Attempt 2 --- lag search added, still FAIL but very differently.}
Correlations roughly doubled, confirming the bug was real and large:

\begin{table}[H]
\centering
\begin{tabular}{lrr}
\toprule
Sequence & Zero-lag (buggy) & With lag search \\
\midrule
0 & 0.42 & \textbf{0.893} \\
2 & --   & 0.721 \\
4 & 0.18 & 0.734 \\
5 & --   & 0.048 \\
6 & 0.26 & 0.097 \\
\bottomrule
\end{tabular}
\caption{Sequence 0 reaches 0.893, just below the pre-registered 0.90 threshold.
The threshold was not lowered to pass the gate.}
\end{table}

\paragraph{Diagnosis: the gate tested the wrong thing.} The candidate list was
built from \emph{standard-catalog colocation} with single-link clustering at
0.5\,km. Standard-catalog absolute errors at Parkfield are 0.2--0.5\,km while
repeaters are $<$20\,m across, so a ``sequence'' built that way mixes different
patches by construction. Low similarity was reporting a broken candidate list,
not absent repeaters --- Parkfield demonstrably has 187 documented sequences over
1987--1998. The field builds repeater catalogs from waveform similarity, not from
locations; this had it backwards.

\subsection{Rebuild on the double-difference catalog --- \accepted}

\textbf{Script:} \texttt{ddrt\_candidates.py}, \texttt{ddrt\_diagnose.py}.

William Ellsworth --- the Ellsworth of Poupinet, Ellsworth \& Fr\'echet (1984),
the paper that introduced doublet monitoring --- supplied the correct catalog:
NCEDC \textbf{Double-Difference relocated} (DDRT), 329 events, 2024-05-01 to
2026-04-01, within 15\,km of SAFOD. DD \emph{relative} locations are accurate to
tens of metres, which is the scale repeaters occupy. His guidance: focus on spots
with multiple dates at one location, prefer larger events, and expect to verify
candidates as ``either repeaters or neighbors.''

Method changes, each addressing a named defect: DD locations; \textbf{complete}
linkage instead of single, so clusters cannot chain; 150\,m diameter instead of
500\,m; magnitude spread $<0.3$; pairs $>30$ days apart.

\paragraph{First run returned zero sequences.} Rather than tune the threshold,
the pair count was decomposed --- which located the real constraint.

\paragraph{A second coverage manifest had been missed.} The coverage list was
built from \texttt{SAFOD\_2024\_2025.csv} alone: 286 days ending 2025-07-25. A
second manifest on the same cable,
\texttt{SAFOD\_vertical\_2026\_01\_28.csv} (182{,}271 rows), adds \textbf{127
further days} and extends coverage to \textbf{2026-01-27}. Union:
\textbf{413 days}, 2024-05-07 to 2026-01-27. Covered events rose 153 $\to$ 215
of 329.

\begin{table}[H]
\centering
\begin{tabular}{rrrrrr}
\toprule
Separation (m) & Pairs & Both covered & $+\Delta M<0.3$ & $+>30$\,d & All filters \\
\midrule
50   & 0    & 0   & 0   & 0    & 0 \\
100  & 14   & 5   & 6   & 10   & 2 \\
150  & 35   & 12  & 11  & 23   & \textbf{3} \\
250  & 116  & 46  & 41  & 84   & \textbf{8} \\
500  & 427  & 186 & 168 & 332  & 51 \\
\bottomrule
\end{tabular}
\caption{Pair-count decomposition after merging both manifests. The binding
constraint is \emph{coverage}, not clustering: tight, magnitude-matched,
well-separated pairs exist, but frequently only one member was recorded.}
\end{table}

\begin{table}[H]
\centering
\begin{tabular}{rrrrrrl}
\toprule
Cluster & $n$ & Diameter & Depth & Mag & km from SAFOD & Separation \\
\midrule
63 & 2 & \textbf{83\,m} & 4.50\,km & \textbf{2.13} & 7.75 & 88 d \\
42 & 2 & 70\,m & 1.76\,km & 0.95 & 2.85 & 39 d \\
\bottomrule
\end{tabular}
\caption{Surviving candidates at 150\,m. Cluster 63 is the strongest target:
M2.13 gives good SNR, 83\,m is within a single M2 rupture dimension, and 88 days
is a usable CWI baseline. G2 on these is queued.}
\end{table}

\subsection{Similarity-first search over all pairs --- \accepted}

\textbf{Scripts:} \texttt{build\_event\_list.py}, \texttt{extract\_all.py} (SLURM
array), \texttt{correlate\_all.py}.

The catalog-first approach was backwards. The field builds repeater catalogs from
waveform similarity and uses location as a check on the answer (REDPy,
Waldhauser--Schaff); I had used location as a filter on the input. Rebuilt: all
207 extractable covered events, all \textbf{21{,}321} pairs, channel-stacked over
channels 100--800, 5--20\,Hz, lag-searched normalised cross-correlation.

The decisive property of doing every pair is that \textbf{the CC distribution is
its own null}. Essentially all pairs are unrelated events, so the bulk of the
histogram measures what ``not a repeater'' looks like on this instrument, in this
band, with this processing --- no threshold need be imported.

\paragraph{A bug that produced a spurious null, found only after the result
looked implausible.} The first run subtracted the across-channel median and then
took the across-channel mean:
\begin{verbatim}
Xb -= np.median(Xb, axis=0)      # remove common mode
s   = Xb[:, i0:i1].mean(axis=0)  # then average channels
\end{verbatim}
mean minus median is approximately zero. Worse, G1's own figure had already shown
the arrival is \emph{flat across all 900 channels} --- near-vertical incidence
means the earthquake \textbf{is} the common mode. Median subtraction is the
standard DAS treatment for common-mode \emph{noise}; here it deleted the signal,
and the reported null was a null for noise. Removing it thickened the tail from 32
to 166 pairs above CC 0.4.

\subsection{HRSN control --- \accepted{} --- REPEATERS CONFIRMED}

\textbf{Script:} \texttt{hrsn\_control.py}.

Even after that fix, no DAS pair reached the CC $\sim$0.9 that defines a
repeater, over 14 months at a site where Schaff \& Waldhauser report $>$40\% of
seismicity is repeating. That inconsistency with the literature could mean either
the processing was blunt or the events genuinely do not repeat --- opposite
remedies, indistinguishable from DAS alone.

HRSN (network BP) is the 13-station borehole array the Parkfield repeater
catalogs were built on, it is public at NCEDC, and it records the same
earthquakes. Correlating the same pairs there, with the processing matched
(identical band, window and lag-searched CC, median across stations):

\begin{table}[H]
\centering
\begin{tabular}{rrrrl}
\toprule
DAS CC & \textbf{HRSN CC} & stations & separation & events \\
\midrule
0.792 & \textbf{0.987} & 8 & 39 d   & 2025-04-02 / 2025-05-11 \\
0.790 & \textbf{0.995} & 5 & 272 d  & 2024-07-08 / 2025-04-06 \\
0.775 & \textbf{0.993} & 3 & 305 d  & 2024-06-23 / 2025-04-23 \\
0.742 & \textbf{0.962} & 8 & 39 d   & 2025-04-02 / 2025-05-11 \\
0.683 & \textbf{0.991} & 5 & 440 d  & 2024-05-13 / 2025-07-27 \\
0.677 & \textbf{0.917} & 5 & 331 d  & 2024-05-10 / 2025-04-06 \\
0.631 & \textbf{0.945} & 4 & 59 d   & 2024-05-10 / 2024-07-08 \\
\bottomrule
\end{tabular}
\caption{Same event pairs, two instruments. \textbf{Seven pairs above CC 0.9 on
HRSN}, with baselines from 39 to 440 days. These are repeating earthquakes.
Random-pair nulls: DAS median 0.124 / max 0.328; HRSN median $-0.004$ / max
0.347.}
\end{table}

Note that 2024-05-10, 2024-07-08 and 2025-04-06 all correlate mutually --- a
three-event sequence spanning 331 days.

\paragraph{Two conclusions, and the second corrects an error of mine.}

\emph{First}, the repeaters exist and have been recorded, with CWI-usable
baselines of 272, 305, 331 and 440 days.

\emph{Second}, the DAS detects every one of them. Against the DAS null --- random
pairs at median 0.124, maximum 0.328 --- confirmed repeaters sit at 0.63--0.79,
two to six times above the null maximum. Every one is an extreme outlier
\emph{in the DAS data}. What DAS does not do is reach 0.9, being single-component
strain rate with 16.34\,m gauge-length averaging against a 250\,Hz borehole
seismometer. \textbf{The earlier ``no repeaters'' conclusion came from applying
HRSN's threshold to DAS} --- exactly the instrument-dependence problem identified
when the 0.9 criterion was first questioned, and then not acted upon.

\subsection{Failures and corrections during this work}

Recorded so that they are not repeated.

\begin{itemize}[leftmargin=*]
  \item \textbf{Timestamp parsing.} \texttt{pd.to\_datetime} without
  \texttt{format='mixed'} rejects the saved candidate times, which vary in
  fractional-second precision. Cost: one job.
  \item \textbf{Empty-frame masking.} Filtering a manifest slice with a boolean
  mask built from \texttt{.map()} on an \emph{empty} frame drops the columns, so
  the subsequent \texttt{sort\_values('t0')} raises \texttt{KeyError} instead of
  returning an empty result. The emptiness check must come first. Present in
  both G1 and G2; patched in both.
  \item \textbf{Buffered output.} \texttt{sbatch --wrap} without \texttt{python
  -u} buffers stdout, so a 47-minute job shows no progress until it exits. Use
  \texttt{python -u}.
  \item \textbf{Zero-lag correlation.} The first G2 estimator omitted the lag
  search, halving every correlation and producing a spurious FAIL. Catalog
  origin-time error alone guarantees this.
  \item \textbf{Wrong catalog.} Candidates were built on standard-catalog
  locations (0.2--0.5\,km error) rather than double-difference relocations
  (tens of metres). The clustering scale must match the location precision.
  \item \textbf{Incomplete coverage list.} Only one of two manifests for the same
  cable was used, undercounting recording days by 127 and truncating coverage six
  months early. Later the same inconsistency appeared again: the merged list was
  used for event SELECTION while extraction still read one manifest, so the second
  member of both candidate pairs was silently skipped.
  \item \textbf{Regex CSV separator.} \texttt{sep=r'\textbackslash s+'} forces
  pandas onto its pure-Python parser; on a 108\,MB manifest that alone exceeded a
  25 minute job limit. Use \texttt{delim\_whitespace} and cache the parsed frame.
  \item \textbf{574k filesystem stats on Lustre.} \texttt{map(os.path.exists)} over
  the whole manifest hammers the metadata server, which this project's
  \texttt{CLAUDE.md} explicitly warns against. Only the few files overlapping each
  event window need checking.
  \item \textbf{Median subtraction destroying the signal.} See the
  similarity-first section. The single most consequential error: it produced a
  believable but meaningless null.
  \item \textbf{Applying an instrument's threshold to a different instrument.}
  The 0.9 criterion is calibrated on seismometers. Used on DAS it hid seven
  genuine repeaters. Thresholds must be calibrated against the null of the
  instrument in hand.
  \item \textbf{Second-resolution cache tags.} Two events in the same second
  collapsed to one cached file, so both ``events'' loaded the identical waveform
  and produced a spurious CC of exactly 1.000.
  \item \textbf{A self-check that checked nothing.} An ``expected under null''
  column recomputed the observed count rather than a model expectation, so it
  trivially agreed with itself.
\end{itemize}

\section{Evidence ladder}

\begin{table}[H]
\centering
\small
\begin{tabular}{p{0.30\linewidth}p{0.15\linewidth}p{0.47\linewidth}}
\toprule
Statement & Status & Basis \\
\midrule
Continuous DAS covers 286 days & \accepted & Manifest filename scan; 392{,}038 rows \\
Catalog events coincide with coverage & \accepted & 305/540 on covered days \\
Candidate sequences exist & \preliminary & 2 at 150\,m, 8 at 250\,m, on DD relocations with 413-day coverage \\
Candidate list is over-merged & \accepted & seq 0 spans 1\,km depth at a 0.5\,km threshold \\
Fiber resolves earthquake coda & \accepted & G1: 785/900 channels, SNR $>3$, 5 events \\
Fiber is degraded below $\sim$850\,m & \accepted & G1 SNR collapse; independently matches Lellouch's 800\,m failure \\
\textbf{Repeating earthquakes present} & \accepted & \textbf{7 pairs at HRSN CC 0.92--0.995, baselines 39--440 d} \\
DAS detects those repeaters & \accepted & DAS CC 0.63--0.79 vs DAS null max 0.328 \\
DAS CC is systematically below HRSN & \accepted & 0.15--0.30 deficit on identical pairs; cause not yet established \\
A three-event sequence exists & \preliminary & 2024-05-10 / 2024-07-08 / 2025-04-06 mutually correlated \\
Pair count is coverage-limited & \accepted & Decomposition: tight matched pairs exist but often only one member recorded \\
$dv/v$ measurable on this fiber & \notrun & Week 2; inputs now exist \\
$dv/v$ resolvable \emph{versus depth} & \notrun & Week 3; the actual deliverable \\
Velocity change confined to top 100\,m & \unresolved & Rubinstein \& Beroza's result; this dataset is the test, not a confirmation \\
\bottomrule
\end{tabular}
\caption{``Accepted'' refers to the stated empirical property only and must not
be silently promoted to a more specific interpretation.}
\end{table}

\section{Next, in order}

\begin{enumerate}[leftmargin=*]
  \item \textbf{G2 completion.} If fewer than 3 sequences confirm, the plan says
  stop --- the candidates were catalog artifacts.
  \item \textbf{Re-cluster} with complete linkage rather than single link, and
  require magnitude spread $<0.3$.
  \item \textbf{Week 2 --- $dv/v$.} Stretching estimator on channel-stacked coda
  first, for maximum SNR, before attempting per-channel work. Three mandatory
  nulls: non-repeating pairs must give no coherent stretch; the two halves of a
  coda must agree; randomised channel order must degrade the stack.
  \item \textbf{Week 3 --- depth profile.} Per-channel $dv/v$ binned to
  $\sim$50\,m, with error bars. This answers Rubinstein \& Beroza.
  \item \textbf{External corroboration.} Cross-check confirmed sequences against
  the Waldhauser--Schaff relocated catalog (\texttt{10.1029/2007JB005479}) and
  Sawi et al.\ (2023), which cover 1984--2019. A sequence active then and now is
  strong independent support.
\end{enumerate}

\section{Reproducibility map}

\begin{table}[H]
\centering
\small
\begin{tabular}{p{0.30\linewidth}p{0.28\linewidth}p{0.34\linewidth}}
\toprule
Script & Output & Purpose \\
\midrule
\texttt{count\_doublets.py} & \texttt{doublet\_candidates.csv} & Catalog intersection \\
\texttt{g1\_coda\_snr.py} & \texttt{g1\_coda\_snr\_seq4.\{npz,png\}} & Gate G1 \\
\texttt{g2\_confirm\_repeaters.py} & \texttt{g2\_confirm\_repeaters.\{npz,png\}}, \texttt{cache/} & Gate G2 \\
\texttt{REPEATERS\_dashboard.ipynb} & figures & Presentation layer \\
\bottomrule
\end{tabular}
\end{table}

Catalog retrieval: USGS FDSN CSV query, box $35.70$--$36.10^\circ$N,
$-120.70$--$-120.20^\circ$E, 2024-05-01 to 2025-10-01, all magnitudes, saved as
\texttt{parkfield\_catalog\_2024\_2025.csv}.

Everything lives in \texttt{notebooks/faultzone/repeaters/}. Nothing in
\texttt{awd\_clean/} is read-modified or written.

\section{Bottom line}

Both preconditions now hold. G1 established that this fiber resolves earthquake
coda across 785 of 900 channels, which nobody had shown. The HRSN control
established that genuine repeating earthquakes were recorded during the DAS
period --- seven pairs at CC 0.92--0.995 with baselines out to 440 days --- and
that the DAS detects all of them as extreme outliers against its own null.

What remains open is not whether the measurement can be made but how to make it
correctly. The DAS returns 0.15--0.30 lower correlation than a co-located
seismometer on identical pairs, and the cause --- single-component strain rate,
16.34\,m gauge-length low-passing, channel stacking, or per-channel SNR --- is not
yet established. A literature methods review is running to settle the processing
chain before the $dv/v$ estimator is built, because the record of this work so far
is that improvised method choices have failed repeatedly and in ways that looked
like results.

\end{document}
